% Arabic abstract (bilingual thesis requirement).
%
% This file is a standalone XeLaTeX document. The main thesis document
% (main.tex) is built with pdflatex + mathptmx (packages.tex) and does not
% support Arabic script, so this abstract is compiled separately with:
%
%     xelatex arabicabstract.tex
%
% and the resulting single-page arabicabstract.pdf
% included in the main document via
% \includepdf in titlepage.tex.
%
% Required packages: fontspec, polyglossia, and the Amiri Arabic font
% (Ubuntu/Debian: texlive-lang-arabic, fonts-hosny-amiri).

\documentclass[12pt,a4paper,oneside]{report}

\usepackage{geometry}
\newgeometry{
    left=2.5cm,
    right=3.0cm,
    top=2.0cm,
    bottom=2.5cm
}

\usepackage{fontspec}
\usepackage{polyglossia}

\setdefaultlanguage{arabic}
\setotherlanguage{english}

\newfontfamily\arabicfont[Script=Arabic]{Amiri}
\usepackage{setspace}
\doublespacing

\pagestyle{empty}

\begin{document}

% ============================================================
% Thesis title
% ============================================================

\begin{center}

{\Large\bfseries
العنوان بالعربي
}\\[0.5cm]

{\large  غرام الشهراني}\\[0.2cm]

{\normalsize
المشرف: أ.د.
} \\[0.7cm]

{\large\bfseries الملخص}


\end{center}

\vspace{0.1cm}

% ============================================================
% Arabic Abstract
% ============================================================

\noindent

ملخص بالعربية لا يتجاوز 700 كلمة  


\end{document}